\section{Conceptual Definitions}

\begin{definition}[Multi-Agent System]\label{def:multi-agent-system}
A multi-agent system consists of independent agents editing the same file.
Each agent records its changes without knowledge of other agents.
\end{definition}

\begin{definition}[Agent]\label{def:agent}
An agent is an entity that edits a file. Each agent has a unique positive integer ID.
\end{definition}

\subsection{The Base File}

\begin{definition}[Base File]\label{def:base-file}
The base file is the original file state before edits. It is a list of text lines.
Example: \code{["name = Alice", "greeting = hello", "print greeting"]}.
\end{definition}

\begin{definition}[Line Number]\label{def:line-number}
Line numbers are 1-indexed and half-open: range \code{[start, stop)} includes \code{start} and excludes \code{stop}.
Example: \code{[2, 4]} covers lines 2 and 3.
\end{definition}

\subsection{Edits and Hunks}

\begin{definition}[Edit]\label{def:edit}
An edit is one of:
\begin{itemize}
  \item Insert: add lines at a position.
  \item Delete: remove lines.
  \item Replace: remove lines and add new ones.
\end{itemize}
\end{definition}

\begin{definition}[Operation]\label{def:operation}
An operation (\code{op}) is one of \code{insert}, \code{delete}, \code{replace}.
\end{definition}

\begin{definition}[Range]\label{def:range}
A range \code{[start, stop)} identifies lines in the base file.
\begin{itemize}
  \item \code{start} included, \code{stop} excluded.
  \item Width = \code{stop - start}.
  \item For insert: \code{start = stop}.
\end{itemize}
Line numbers are 1-based.
\end{definition}

\begin{definition}[Before Lines]\label{def:before-lines}
\code{before} holds the exact lines from the base file at \code{range}. For insert it is empty.
\end{definition}

\begin{definition}[After Lines]\label{def:after-lines}
\code{after} holds the lines to insert. For delete it is empty.
\end{definition}

\begin{definition}[Hunk]\label{def:hunk}
A hunk is one atomic edit. It contains \code{seq}, \code{op}, \code{range}, \code{before}, \code{after}.
See Definitions~\ref{def:range}, \ref{def:operation}, \ref{def:before-lines}, \ref{def:after-lines}.
\end{definition}

\subsection{Metadata}

\begin{definition}[Agent ID]\label{def:agent-id}
Agent ID is a unique positive integer $\geq 1$. It appears at top level of every diff.
\end{definition}

\begin{definition}[Timestamp]\label{def:timestamp}
Timestamp is an ISO 8601 date-time string.
\end{definition}

\begin{definition}[File Path]\label{def:file-path}
File path is a non-empty string identifying the target file.
\end{definition}

\subsection{Diff}

\begin{definition}[Base Line Count]\label{def:base-line-count}
Base line count is the number of lines in the base file ($\geq 0$).
\end{definition}

\begin{definition}[Diff]\label{def:diff}
A diff contains agent ID, file path, timestamp, base line count, and ordered list of hunks.
\end{definition}

\subsection{Constraints}

\begin{definition}[Non-Overlap]\label{def:non-overlap}
No two hunks in one diff may have overlapping ranges: \code{a < d \land c < b} for ranges \code{[a,b)}, \code{[c,d)}.
\end{definition}

\begin{definition}[Schema Validity]\label{def:schema-validity}
A diff satisfies the JSON schema if all required fields, types, bounds, shapes, and non-overlap hold.
\end{definition}

\begin{definition}[Base Alignment]\label{def:base-alignment}
A hunk is base-aligned when \code{before} matches the base file slice at its range.
A diff is base-aligned when all hunks are aligned and \code{base_line_count} matches base length.
\end{definition}
